\documentclass[11pt,a4paper]{article}

\usepackage[T1]{fontenc}
\usepackage[utf8]{inputenc}
\usepackage[french]{babel}
\usepackage{lmodern}
\usepackage{microtype}
\usepackage{geometry}
\usepackage{amsmath,amssymb,mathtools}
\usepackage{booktabs}
\usepackage{enumitem}
\usepackage{xcolor}
\usepackage{titlesec}
\usepackage{fancyhdr}
\usepackage{hyperref}
\usepackage[most]{tcolorbox}
\usepackage{lastpage}

\geometry{top=22mm,bottom=24mm,left=24mm,right=24mm,headheight=14pt}
\definecolor{ink}{HTML}{172033}
\definecolor{accent}{HTML}{315A8A}
\definecolor{accentlight}{HTML}{EAF1F8}
\definecolor{rulegray}{HTML}{D7DEE8}
\definecolor{muted}{HTML}{5C6678}

\hypersetup{
  colorlinks=true,
  linkcolor=accent,
  urlcolor=accent,
  citecolor=accent,
  pdftitle={Symétries de demi-période pour les suites modulo 3},
  pdfsubject={Formalisation mathématique et audit algorithmique},
  pdfkeywords={récurrence, modulo 3, antipériodicité, Fourier, matrice compagnon}
}

\pagestyle{fancy}
\fancyhf{}
\fancyhead[L]{\small\color{muted}Symétries de demi-période modulo 3}
\fancyhead[R]{\small\color{muted}Formalisation et audit}
\fancyfoot[C]{\small\color{muted}\thepage\,/\,\pageref*{LastPage}}
\renewcommand{\headrulewidth}{0.3pt}
\renewcommand{\footrulewidth}{0pt}

\titleformat{\section}
  {\Large\bfseries\color{ink}}
  {\thesection.}{0.6em}{}
  [\vspace{0.2em}\color{rulegray}\titlerule]
\titleformat{\subsection}
  {\large\bfseries\color{accent}}
  {\thesubsection}{0.6em}{}
\titlespacing*{\section}{0pt}{2.2ex plus .5ex minus .2ex}{1.3ex}
\titlespacing*{\subsection}{0pt}{1.7ex plus .4ex minus .2ex}{0.8ex}

\setlist[itemize]{leftmargin=1.6em,itemsep=0.25em,topsep=0.4em}
\setlist[enumerate]{leftmargin=1.8em,itemsep=0.25em,topsep=0.4em}
\setlength{\parindent}{0pt}
\setlength{\parskip}{0.55em}
\allowdisplaybreaks

\newtcolorbox{summarybox}{
  colback=accentlight,
  colframe=accent,
  boxrule=0.6pt,
  arc=2mm,
  left=4mm,right=4mm,top=3mm,bottom=3mm,
  before skip=1em,after skip=1em
}
\newcommand{\F}{\mathbb F}
\newcommand{\Z}{\mathbb Z}
\newcommand{\concat}{\mathbin{\Vert}}

\begin{document}

\begin{titlepage}
\thispagestyle{empty}
\vspace*{22mm}
{\color{accent}\rule{\textwidth}{1.2pt}}\\[15mm]
{\fontsize{27}{33}\selectfont\bfseries\color{ink}
Symétries de demi-période\\[2mm]
pour les suites modulo 3\par}
\vspace{9mm}
{\Large\color{accent}Formalisation mathématique et audit algorithmique\par}
\vspace{13mm}
\begin{summarybox}
Cette note distingue rigoureusement trois phénomènes souvent confondus :
la négation de cycles, l'antipériodicité de demi-période et le renversement.
Elle établit le décalage opposé unique d'un mot primitif non nul, donne des
caractérisations combinatoire, spectrale et matricielle, puis relie ces
résultats à un analyseur exhaustif de récurrences modulaires.
\end{summarybox}
\vfill
{\large\color{muted}Dépôt : \texttt{Modular-recurrence-symetry-analyzer-}\par}
{\large\color{muted}Version PDF consolidée -- 5 août 2026\par}
\vspace{12mm}
{\color{accent}\rule{\textwidth}{1.2pt}}
\end{titlepage}

\tableofcontents
\clearpage

\section{Objet}

Cette note formalise les phénomènes décrits dans le dépôt
\href{https://github.com/59200/k-bonacci-pisano-finder}{\texttt{59200/k-bonacci-pisano-finder}}
sous les expressions informelles « périodes complémentaires à trois » et
« auto-complémentaires lorsqu'elles sont coupées en deux ».

Le cadre invariant est celui des mots périodiques sur le corps
\[
\F_3=\{0,1,2\},
\]
où l'involution additive est
\[
\nu(x)=-x\pmod 3,
\qquad
\nu(0)=0,\quad \nu(1)=2,\quad \nu(2)=1.
\]

L'expression « complémentaire à 3 » doit donc être remplacée par
\textbf{opposé additif modulo 3} ou \textbf{complément par négation modulo 3}.
La somme décimale des représentants n'est pas toujours égale à 3, puisque
0 est envoyé sur 0.

\section{Trois propriétés distinctes}

Soit \(w=(w_0,\ldots,w_{T-1})\in\F_3^T\) un mot de période primitive \(T\).

\subsection{Paire de périodes opposées}

Deux périodes \(w\) et \(v\) sont une paire opposée si
\[
v_n=-w_n\pmod 3
\]
pour tout \(n\). Elles peuvent appartenir à deux cycles distincts.

\subsection{Antipériodicité de demi-période}

Le mot \(w\) est \textbf{antipériodique à demi-période} si \(T=2h\) et
\[
w_{n+h}=-w_n\pmod 3
\]
pour tout \(n\). Il s'écrit alors
\[
w=H\concat(-H),
\qquad
H=(w_0,\ldots,w_{h-1}).
\]

Exemples :
\[
01120221=0112\concat0221,
\]
\[
011022=011\concat022,
\]
\[
01220211=0122\concat0211.
\]

\subsection{Renversement de la seconde moitié}

L'opération
\[
\mathcal R(a_0,\ldots,a_{h-1})=(a_{h-1},\ldots,a_0)
\]
est indépendante de la négation. Dans l'exemple de Fibonacci,
\[
\mathcal R(0221)=1220.
\]
Ainsi, \(1220\) est le \textbf{mot antipodal renversé} de \(0112\), et non
la seconde moitié elle-même.

\section{Théorème de l'unique décalage opposé}

\begin{summarybox}
\textbf{Théorème.} Soit \(w\) un mot primitif non nul de période \(T\).
S'il existe un décalage \(r\) tel que
\[
w_{n+r}=-w_n
\]
pour tout \(n\), alors \(T\) est pair et
\[
r\equiv \frac{T}{2}\pmod T.
\]
\end{summarybox}

\subsection*{Preuve}

En appliquant deux fois le décalage,
\[
w_{n+2r}=w_n.
\]
La primitivité implique \(T\mid 2r\). Le décalage \(r\) n'est pas nul
modulo \(T\), car un mot non nul sur \(\F_3\) ne satisfait pas \(w=-w\).
L'élément non trivial d'ordre 2 du groupe cyclique \(\Z/T\Z\) existe
uniquement lorsque \(T\) est pair et vaut \(T/2\). \hfill\(\square\)

\subsection*{Conséquence}

Une période primitive non nulle est soit :
\begin{enumerate}
\item envoyée par la négation sur un cycle distinct ;
\item invariante par négation, auquel cas elle est nécessairement
antipériodique à demi-période.
\end{enumerate}

Cette dichotomie corrige la confusion entre les couples
\[
00111201,\quad 00222102
\]
et les périodes intrinsèquement antipériodiques comme \(011022\).

\section{Conséquences combinatoires}

Si
\[
w=H\concat(-H),
\]
alors
\[
\#\{n:w_n=1\}=\#\{n:w_n=2\},
\]
\[
\#\{n:w_n=0\}\equiv0\pmod2,
\]
et
\[
\sum_{n=0}^{T-1}w_n=0
\quad\text{dans }\F_3.
\]
Ces conditions sont nécessaires, mais non suffisantes.

En utilisant le relèvement centré
\[
\widetilde 0=0,\qquad
\widetilde 1=1,\qquad
\widetilde 2=-1,
\]
le polynôme générateur s'écrit
\[
W(X)
=\sum_{n=0}^{2h-1}\widetilde w_nX^n
=(1-X^h)\sum_{n=0}^{h-1}\widetilde w_nX^n.
\]

\section{Caractérisation de Fourier}

Soit
\[
\widehat w(k)
=\sum_{n=0}^{2h-1}
\widetilde w_n e^{-2\pi i kn/(2h)}.
\]
Alors
\[
w_{n+h}=-w_n
\]
si et seulement si
\[
\widehat w(k)=0
\]
pour tout indice pair \(k\).

\subsection*{Preuve directe}

Lorsque \(w_{n+h}=-w_n\),
\[
\widehat w(k)
=\left(1-(-1)^k\right)
\sum_{n=0}^{h-1}
\widetilde w_n e^{-2\pi i kn/(2h)}.
\]
Le facteur est nul pour \(k\) pair.

Réciproquement, si tous les modes pairs sont nuls, le mot appartient au
sous-espace engendré par les modes impairs. Le décalage de \(h\) agit sur
le mode \(k\) par \((-1)^k=-1\), donc il agit par \(-I\) sur le mot.

Cette caractérisation donne un diagnostic spectral exact, distinct d'une
simple comparaison visuelle des deux moitiés.

\section{Suites linéaires modulo 3}

Considérons une récurrence d'ordre \(k\)
\[
u_{n+k}
=c_0u_n+c_1u_{n+1}+\cdots+c_{k-1}u_{n+k-1}
\pmod3.
\]
Son état est
\[
s_n=(u_n,\ldots,u_{n+k-1})^{\mathsf T}
\]
et son évolution est
\[
s_{n+1}=Ms_n,
\]
où \(M\) est la matrice compagnon.

\subsection{Périodicité pure}

Le déterminant de \(M\) vaut, au signe près, \(c_0\). Sur \(\F_3\),
\[
M\text{ est inversible}
\iff c_0\neq0.
\]
Dans ce cas, toute orbite d'état est purement périodique. Si \(c_0=0\),
une prépériode peut apparaître et doit être distinguée du cycle.

\subsection{Négation des cycles}

La linéarité donne
\[
M(-s)=-M(s).
\]
La négation envoie donc tout cycle sur un cycle de même longueur. Pour un
cycle non nul, deux cas seulement sont possibles :
\begin{itemize}
\item deux cycles distincts échangés par négation ;
\item un même cycle invariant, nécessairement antipériodique à demi-période.
\end{itemize}
Le cycle nul est l'unique exception dégénérée : il est fixé point par point
par la négation et ne définit pas une antipériode primitive non triviale.

\subsection{Critère sur une graine}

Pour une graine \(s_0\), la relation
\[
M^hs_0=-s_0
\]
implique, pour toute sortie linéaire de l'état,
\[
u_{n+h}=-u_n.
\]
Lorsque \(s_0\neq0\) et que le cycle est primitif, cette relation fournit
l'antipériode de demi-période décrite plus haut.

\subsection{Critère uniforme}

Toutes les graines possèdent la même relation d'antipériode \(h\) si et
seulement si
\[
M^h=-I.
\]
Équivalemment, le polynôme minimal \(\mu_M\) vérifie
\[
\mu_M(X)\mid X^h+1
\quad\text{dans }\F_3[X].
\]
Il en résulte
\[
M^{2h}=I.
\]

\section{Exemple de Fibonacci}

Pour
\[
M=
\begin{pmatrix}
0&1\\
1&1
\end{pmatrix}
\quad\text{sur }\F_3,
\]
on calcule
\[
M^4=-I,
\qquad
M^8=I.
\]
La période issue de la graine \((0,1)\) est donc
\[
01120221=0112\concat(-0112).
\]
Sa seconde moitié renversée est
\[
\mathcal R(0221)=1220.
\]
La coïncidence numérique, en lecture décimale,
\[
1220=L_{14}+F_{14}=2F_{15}
\]
est une identité d'encodage supplémentaire. Elle n'est pas une conséquence
générale de l'antipériodicité.

\section{Généralisation aux suites arbitraires modulo 3}

Aucune hypothèse de récurrence n'est nécessaire pour analyser un mot
périodique donné. Pour toute période primitive, on peut calculer :
\begin{itemize}
\item son cycle sous les rotations ;
\item son opposé additif ;
\item ses symétries affines \(w_{n+r}=aw_n+b\) ;
\item ses symétries de renversement \(w_{r-n}=aw_n+b\) ;
\item son éventuelle antipériodicité ;
\item ses modes de Fourier pairs et impairs.
\end{itemize}
La théorie matricielle intervient uniquement lorsqu'une récurrence linéaire
déclarée est disponible.

\section{Extension au module général}

Sur \(\Z/m\Z\), l'involution naturelle reste
\[
x\longmapsto-x.
\]
La définition
\[
w_{n+h}=-w_n\pmod m
\]
reste valable. Pour \(m=2\), la négation est l'identité et la notion
d'antipériode se confond avec une période ordinaire. Pour \(m>2\), la
distinction demeure non triviale.

Le terme « complément à \(m\) » doit être évité comme définition
algébrique, car il dépend du choix des représentants entiers. Le terme
invariant est « opposé additif modulo \(m\) ».

\section{Audit algorithmique}

Une période d'une récurrence d'ordre \(k\) modulo \(m\) doit être détectée
sur l'espace fini des états
\[
(\Z/m\Z)^k,
\]
et non par recherche d'un bloc répété dans un préfixe de longueur arbitraire.

Le script fourni utilise :
\begin{itemize}
\item le codage compact des états en base \(m\) ;
\item l'algorithme de Brent pour une graine,
\[
O(\mu+\lambda)\text{ temps},\qquad O(1)\text{ mémoire};
\]
\item une décomposition exacte du graphe fonctionnel pour toutes les graines,
\[
O(m^k)\text{ temps};
\]
\item la parallélisation entre familles de coefficients ;
\item une classification séparée des cycles opposés et des cycles
antipériodiques ;
\item un contrôle matriciel de \(M^h=-I\).
\end{itemize}

Aucun noyau Mathematica ni wrapper externe n'est requis.

\section{Résultats reproductibles du balayage}

Le balayage des quinze familles nommées dans le dépôt retrouve notamment,
pour la famille tritetranacci, les classes suivantes :
\begin{itemize}
\item paires opposées distinctes de longueurs \(24\) et \(8\) ;
\item périodes antipériodiques \(011022\), \(01220211\) et \(12\).
\end{itemize}

Un second balayage porte sur les \(119\) vecteurs de coefficients binaires
non nuls d'ordres \(2\) à \(6\) :
\[
\sum_{k=2}^{6}(2^k-1)=119.
\]
Les calculs produits par le script trouvent :
\begin{itemize}
\item \(13\) familles satisfaisant un critère global \(M^h=-I\) ;
\item \(88\) familles possédant au moins un cycle antipériodique non trivial ;
\item \(96\) familles possédant au moins une paire de cycles opposés distincts.
\end{itemize}

Ces nombres sont des résultats expérimentaux exhaustifs dans cette fenêtre
finie. Ils sont reproductibles à partir des fichiers
\texttt{mod3\_symmetry\_scan\_summary.json},
\texttt{binary\_recurrence\_families\_mod3.json} et
\texttt{binary\_recurrence\_families\_mod3\_summary.csv} fournis dans le dépôt.

\section*{Conclusion}
\addcontentsline{toc}{section}{Conclusion}

La négation modulo 3 agit comme une involution structurante sur les mots
périodiques et les cycles de récurrences linéaires. Pour un mot primitif non
nul, son invariance par cette involution ne peut se produire que par le
décalage unique de demi-période. Cette propriété se reconnaît de façon
équivalente par la factorisation du polynôme générateur, l'annulation des
modes de Fourier pairs ou, dans le cadre linéaire, par une relation matricielle
\(M^h=-I\). L'analyse algorithmique sépare ainsi proprement les symétries
intrinsèques des cycles et les simples associations entre cycles opposés.

\end{document}
